\section{Présentation du corpus d'étude}

\subsection{Choix du corpus}
Cette étude porte sur un corpus composé exclusivement de romans d'Émile Zola. Le choix a été fait de ne pas inclure les nouvelles, les textes journalistiques ni les essais, afin de conserver un ensemble relativement homogène du point de vue générique. Les romans offrent en effet une matière textuelle plus longue et plus développée, ce qui permet d'observer plus finement les thèmes récurrents, les milieux sociaux et les motifs narratifs présents dans l'œuvre de Zola.

L'exclusion des nouvelles répond également à une contrainte méthodologique. Ces textes sont généralement plus courts que les romans et leur intégration aurait pu créer un déséquilibre dans le corpus. Dans le cadre d'une analyse thématique fondée sur des segments textuels, les différences importantes de longueur entre les œuvres peuvent influencer la représentation des thèmes. Limiter le corpus aux romans permet donc de réduire cet effet et de travailler sur un ensemble plus comparable.

Le corpus comprend les œuvres suivantes :

\begin{longtable}{p{1.6cm}p{4.2cm}p{8cm}}
\caption{Corpus des œuvres de Zola utilisées dans l'analyse}\\
\toprule
\textbf{Date} & \textbf{Titre} & \textbf{Fichier} \\
\midrule
\endfirsthead

\toprule
\textbf{Date} & \textbf{Titre} & \textbf{Fichier} \\
\midrule
\endhead

1865 & \textit{La confession de Claude} & \texttt{1865\_La\_confession\_de\_Claude.\_clean.txt} \\
1866 & \textit{Le vœu d'une morte} & \texttt{1866\_Le\_voeu\_d\_une\_morte.\_clean.txt} \\
1867 & \textit{Les mystères de Marseille} & \texttt{1867\_Les\_mysteres\_de\_Marseille.\_clean.txt} \\
1867 & \textit{Thérèse Raquin} & \texttt{1867\_Therese\_Raquin.\_clean.txt} \\
1868 & \textit{Madeleine Férat} & \texttt{1868\_Madeleine\_Ferat.\_clean.txt} \\
1871 & \textit{La fortune des Rougon} & \texttt{1871\_1\_La\_fortune\_des\_Rougon.\_clean.txt} \\
1872 & \textit{La curée} & \texttt{1871\_2\_La\_curee.\_clean.txt} \\
1873 & \textit{Le ventre de Paris} & \texttt{1873\_3\_Le\_ventre\_de\_Paris.\_clean.txt} \\
1874 & \textit{La conquête de Plassans} & \texttt{1874\_4\_La\_conquete\_de\_Plassans.\_clean.txt} \\
1875 & \textit{La faute de l'abbé Mouret} & \texttt{1875\_5\_La\_faute\_de\_l\_abbe\_Mouret.\_clean.txt} \\
1876 & \textit{Son Excellence Eugène Rougon} & \texttt{1876\_6\_Son\_Excellence\_Eugene\_Rougon.\_clean.txt} \\
1877 & \textit{L'assommoir} & \texttt{1877\_7\_L\_assommoir.\_clean.txt} \\
1878 & \textit{Une page d'amour} & \texttt{1878\_8\_Une\_page\_d\_amour.\_clean.txt} \\
1880 & \textit{Nana} & \texttt{1880\_9\_Nana.\_clean.txt} \\
1882 & \textit{Pot-Bouille} & \texttt{1882\_10\_Pot-bouille.\_clean.txt} \\
1883 & \textit{Au Bonheur des dames} & \texttt{1883\_11\_Au\_Bonheur\_des\_dames.\_clean.txt} \\
1884 & \textit{La joie de vivre} & \texttt{1884\_12\_La\_joie\_de\_vivre.\_clean.txt} \\
1885 & \textit{Germinal} & \texttt{1885\_13\_Germinal.\_clean.txt} \\
1886 & \textit{L'œuvre} & \texttt{1886\_14\_L\_oeuvre.\_clean.txt} \\
1887 & \textit{La terre} & \texttt{1887\_15\_La\_terre.\_clean.txt} \\
1888 & \textit{Le rêve} & \texttt{1888\_16\_Le\_reve.\_clean.txt} \\
1890 & \textit{La bête humaine} & \texttt{1890\_17\_La\_bete\_humaine.\_clean.txt} \\
1891 & \textit{L'argent} & \texttt{1891\_18\_L\_argent.\_clean.txt} \\
1892 & \textit{La débâcle} & \texttt{1892\_19\_La\_debacle.\_clean.txt} \\
1893 & \textit{Le docteur Pascal} & \texttt{1893\_20\_Le\_docteur\_Pascal.\_clean.txt} \\
1894 & \textit{Lourdes} & \texttt{1894\_1\_Lourdes.\_clean.txt} \\
1896 & \textit{Rome} & \texttt{1896\_2\_Rome.\_clean.txt} \\
1898 & \textit{Paris} & \texttt{1898\_3\_Paris.\_clean.txt} \\
1899 & \textit{Fécondité} & \texttt{1899\_1\_Fecondite.\_clean.txt} \\
1901 & \textit{Travail} & \texttt{1901\_2\_Travail.\_clean.txt} \\
1903 & \textit{Vérité} & \texttt{1903\_3\_Verite.\_clean.txt} \\

\bottomrule
\end{longtable}

Le corpus est composé de 31 romans, qui couvrent les principales périodes de la production romanesque de Zola. Il comprend :

\begin{itemize}
    \item 5 romans antérieurs au cycle des \textit{Rougon-Macquart} ;
    \item 20 romans appartenant au cycle des \textit{Rougon-Macquart} ;
    \item 3 romans du cycle des \textit{Trois Villes} ;
    \item 3 romans publiés du cycle des \textit{Quatre Évangiles}.
\end{itemize}

Pour ce qui est des \textit{Quatre Évangiles}, \textit{Vérité} paraît à titre posthume. \textit{Justice} ne paraîtra jamais, l'ouvrage étant resté à l'état d'ébauche au moment de la mort de l'écrivain.

Ce corpus permet donc de couvrir une large période de la carrière de Zola, depuis ses premiers romans jusqu'à ses dernières œuvres. Cette amplitude chronologique est importante pour l'analyse, car elle permet d'observer non seulement les thèmes récurrents de l'œuvre romanesque, mais aussi leur possible évolution au fil du temps\footnote{En vue d'une éventuelle analyse ultérieure par dynamic topic modeling sur ce même corpus.}.


\subsection{Préparation des données brutes en données propres}

Les romans au format \texttt{.txt} correspondent à des versions nettoyées, c'est-à-dire à des textes dont les éléments non textuels ont été supprimés. Ces éléments comprennent notamment les titres, les chapitres, les notes de bas de page et d'autres informations éditoriales qui pourraient perturber l'analyse. Ce nettoyage permet d'obtenir une représentation plus précise des thèmes présents dans les romans.

De plus, le choix a été fait de conserver une seule phrase par ligne. Ce format facilite l'analyse textuelle, notamment pour les méthodes de la NMF et de K-means, qui nécessitent une structure de données homogène pour fonctionner correctement.

Cette logique de formatage a également été appliquée en prévision des analyses futures. Les romans issus de prochains traitements pourront ainsi être intégrés de manière cohérente au corpus. Ce choix permet de limiter les différences de structure entre les fichiers et d'obtenir un format homogène pour l'analyse sans trop complexifier le processus de préparation des données.

\begin{longtable}{p{8cm}r}
\caption{Nombre de tokens par roman dans l'ordre chronologique}  \\
\toprule
\textbf{Titre du roman} & \textbf{Nombre de tokens} \\
\midrule
\endfirsthead

\toprule
\textbf{Titre du roman} & \textbf{Nombre de tokens} \\
\midrule
\endhead

\textit{La confession de Claude} & 47253 \\
\textit{Le vœu d'une morte} & 38722 \\
\textit{Les mystères de Marseille} & 125253 \\
\textit{Thérèse Raquin} & 66363 \\
\textit{Madeleine Férat} & 97318 \\
\textit{La fortune des Rougon} & 116685 \\
\textit{La curée} & 105120 \\
\textit{Le ventre de Paris} & 110442 \\
\textit{La conquête de Plassans} & 113245 \\
\textit{La faute de l'abbé Mouret} & 114667 \\
\textit{Son Excellence Eugène Rougon} & 126186 \\
\textit{L'assommoir} & 158340 \\
\textit{Une page d'amour} & 102185 \\
\textit{Nana} & 141573 \\
\textit{Pot-Bouille} & 134844 \\
\textit{Au Bonheur des dames} & 147521 \\
\textit{La joie de vivre} & 117255 \\
\textit{Germinal} & 164473 \\
\textit{L'œuvre} & 130613 \\
\textit{La terre} & 163687 \\
\textit{Le rêve} & 66379 \\
\textit{La bête humaine} & 125418 \\
\textit{L'argent} & 143010 \\
\textit{La débâcle} & 187046 \\
\textit{Le docteur Pascal} & 112713 \\
\textit{Lourdes} & 173511 \\
\textit{Rome} & 233958 \\
\textit{Paris} & 177389 \\
\textit{Fécondité} & 220131 \\
\textit{Travail} & 198712 \\
\textit{Vérité} & 223274 \\

\bottomrule
\end{longtable}

Le tableau fournit une description quantitative du corpus en indiquant, pour chaque roman, le nombre de tokens obtenus après nettoyage. Les œuvres sont présentées dans l'ordre chronologique, ce qui permet d'observer l'évolution du volume textuel au sein de la production romanesque de Zola. Les écarts sont importants entre les romans courts du début de carrière, comme \textit{Le vœu d'une morte}, et les œuvres plus volumineuses de la fin de carrière, notamment \textit{Rome}, \textit{Fécondité}, \textit{Travail} et \textit{Vérité}.

Cette information est importante pour l'interprétation des résultats, car la longueur des textes peut avoir un effet sur la représentation des thèmes. Un roman plus long fournit mécaniquement davantage d'occurrences lexicales et davantage de segments analysables. La segmentation en paquets de phrases vise donc à produire des unités d'analyse plus comparables, tout en conservant suffisamment de contexte pour identifier les thèmes dominants.
\subsection{Préparation du corpus pour l'analyse thématique}

Pour préparer le corpus à l'analyse thématique, nous avons adopté une procédure visant à produire des unités textuelles comparables et exploitables par les modèles.

Les romans ont été segmentés en paquets de phrases\footnote{Le code pour segmenter les romans est disponible en annexe, dans la section « Segmentation du corpus ».}. Chaque paquet est composé de 40 phrases. Cette taille a été retenue car elle offre un compromis entre deux exigences : conserver suffisamment de contexte pour identifier les thèmes présents dans les romans, tout en permettant une analyse plus fine qu'un découpage à l'échelle du chapitre ou de l'œuvre entière.

Ce choix reste toutefois méthodologique et partiellement arbitraire. Plusieurs tailles de paquets ont été testées\footnote{Différents essais ont été réalisés avec des paquets de 10, 50, 70 et 100 phrases.}. Les paquets de 40 phrases ont été retenus, car ils produisaient les résultats les plus cohérents et les plus interprétables parmi les configurations testées. Ce choix n'est donc pas parfait, mais il constitue le meilleur compromis observé dans le cadre de cette analyse.

D'autres modes de segmentation ont également été envisagés, notamment un découpage par chapitre ou par paragraphe. Le découpage par chapitre a été écarté, car les chapitres ne constituent pas des unités suffisamment homogènes en longueur et en contenu. Certains chapitres sont très courts, tandis que d'autres sont beaucoup plus longs, ce qui rend la comparaison entre unités d'analyse plus difficile.

Le découpage par paragraphe a également été étudié, mais il posait plusieurs problèmes. D'abord, les paragraphes varient fortement en longueur : certains sont très courts, tandis que d'autres sont beaucoup plus développés. Cette difficulté pourrait être partiellement corrigée par un filtrage des paragraphes trop courts ou trop longs, mais cela introduirait une étape supplémentaire et rendrait le traitement moins homogène.

Ensuite, l'architecture matérielle des textes n'est pas toujours conservée de manière fiable dans les fichiers disponibles. Certains romans ne présentent plus une mise en page avec des paragraphes clairement identifiables, en particulier lorsque les textes proviennent de formats différents, comme des fichiers \texttt{.txt}, \texttt{.pdf} ou \texttt{.epub}. Or, ce travail s'inscrit dans une perspective d'élargissement du corpus à d'autres romans et à d'autres sources textuelles. Dans ce contexte, la segmentation par phrase apparaît plus facilement généralisable à des formats hétérogènes.

Enfin, la structure interne des romans pose un dernier problème. Les dialogues peuvent être considérés comme des paragraphes indépendants, ce qui produit parfois des unités d'analyse très courtes, composées de seulement quelques mots. Là encore, un filtrage serait possible, mais il complexifierait le protocole et risquerait d'introduire des différences de traitement entre les textes.

Pour ces raisons, le découpage en paquets de 40 phrases a été retenu. Il permet de produire des unités d'analyse relativement homogènes, applicables à l'ensemble du corpus, tout en conservant une quantité suffisante de contexte pour l'analyse thématique.
\newpage
\subsection{Statistiques descriptives du corpus}

\begin{longtable}{p{8cm}r}
\caption{Nombre de paquets de phrases par roman} \\
\toprule
\textbf{Titre du roman} & \textbf{Nombre de paquets} \\
\midrule
\endfirsthead

\toprule
\textbf{Titre du roman} & \textbf{Nombre de paquets} \\
\midrule
\endhead

\textit{La confession de Claude} & 65 \\
\textit{Le vœu d'une morte} & 60 \\
\textit{Les mystères de Marseille} & 184 \\
\textit{Thérèse Raquin} & 80 \\
\textit{Madeleine Férat} & 124 \\
\textit{La fortune des Rougon} & 154 \\
\textit{La curée} & 123 \\
\textit{Le ventre de Paris} & 136 \\
\textit{La conquête de Plassans} & 165 \\
\textit{La faute de l'abbé Mouret} & 171 \\
\textit{Son Excellence Eugène Rougon} & 191 \\
\textit{L'assommoir} & 228 \\
\textit{Une page d'amour} & 166 \\
\textit{Nana} & 225 \\
\textit{Pot-Bouille} & 207 \\
\textit{Au Bonheur des dames} & 198 \\
\textit{La joie de vivre} & 164 \\
\textit{Germinal} & 225 \\
\textit{L'œuvre} & 163 \\
\textit{La terre} & 222 \\
\textit{Le rêve} & 82 \\
\textit{La bête humaine} & 154 \\
\textit{L'argent} & 155 \\
\textit{La débâcle} & 221 \\
\textit{Le docteur Pascal} & 135 \\
\textit{Lourdes} & 198 \\
\textit{Rome} & 231 \\
\textit{Paris} & 206 \\
\textit{Fécondité} & 259 \\
\textit{Travail} & 208 \\
\textit{Vérité} & 227 \\

\bottomrule
\end{longtable}

Ce tableau indique le nombre de paquets de phrases obtenus pour chaque roman après segmentation du corpus. Les écarts observés reflètent principalement les différences de longueur entre les œuvres : les textes les plus courts, comme \textit{Le vœu d'une morte} ou \textit{La confession de Claude}, produisent peu de paquets, tandis que les romans plus volumineux, comme \textit{Fécondité}, \textit{Rome} ou \textit{L'assommoir}, en produisent davantage.

\subsubsection{Analyse des statistiques descriptives du nombre de paquets par roman}
\begin{table}[H]
\centering
\caption{Statistiques descriptives du nombre de paquets par roman}
\begin{tabular}{lr}
\toprule
\textbf{Indicateur} & \textbf{Valeur} \\
\midrule
Nombre de romans & 31 \\
Moyenne & 171.84 \\
Écart-type & 52.40 \\
Minimum & 60 \\
Premier quartile & 145 \\
Médiane & 171 \\
Troisième quartile & 214.5 \\
Maximum & 259 \\
\bottomrule
\end{tabular}
\end{table}


Les statistiques descriptives du nombre de paquets par roman permettent d'évaluer la structure quantitative du corpus après segmentation. Le corpus comprend 31 romans, avec une moyenne d'environ 172 paquets de phrases par œuvre. La médiane, située à 171 paquets, est très proche de cette moyenne. Cette proximité indique que la distribution générale du nombre de paquets n'est pas fortement asymétrique : la majorité des romans se situe autour d'un volume intermédiaire, même si certains textes s'écartent nettement de cette tendance.

L'écart-type, d'environ 52 paquets, montre toutefois une dispersion importante autour de la moyenne. Cette variation s'explique principalement par les différences de longueur entre les romans. Les œuvres les plus courtes, comme \textit{Le vœu d'une morte}, produisent seulement 60 paquets, tandis que les textes les plus longs, comme \textit{Fécondité}, atteignent 259 paquets. L'écart entre le minimum et le maximum est donc notable, puisqu'il existe presque un rapport de un à quatre entre le roman le moins segmenté et le roman le plus segmenté.

Les quartiles permettent de préciser cette répartition. Le premier quartile, situé à 145 paquets, indique qu'un quart des romans contient 145 paquets ou moins. Le troisième quartile, situé à 214,5 paquets, indique qu'environ trois quarts des romans restent sous ce seuil. La moitié centrale du corpus se situe donc entre 145 et 214,5 paquets, ce qui montre que la majorité des œuvres conserve une taille relativement comparable après segmentation. Les valeurs extrêmes ne remettent donc pas en cause l'homogénéité générale du corpus, mais elles doivent être prises en compte dans l'interprétation des résultats.

Le coefficient de variation est d'environ 30 \%, puisqu'il correspond au rapport entre l'écart-type et la moyenne. Cette valeur signale une variabilité modérée : les romans ne sont pas tous de taille équivalente, mais les écarts restent attendus pour un corpus littéraire composé d'œuvres publiées à des périodes différentes et appartenant à plusieurs cycles romanesques. Cette dispersion est donc méthodologiquement acceptable, à condition de garder à l'esprit que les romans les plus longs fournissent davantage d'unités d'analyse.

Cette information est importante pour l'analyse thématique. Chaque paquet de phrases constitue une unité soumise aux méthodes de topic modeling. Par conséquent, un roman plus long contribue mécaniquement à un plus grand nombre d'observations qu'un roman court. Les thèmes présents dans les romans volumineux peuvent donc être davantage représentés dans les résultats globaux. La segmentation en paquets de 40 phrases permet néanmoins de réduire cet effet, car elle transforme les romans en unités textuelles comparables tout en conservant suffisamment de contexte pour identifier des thèmes cohérents.

Ainsi, ces statistiques montrent que le corpus présente un équilibre satisfaisant pour une analyse thématique. Les différences de taille entre les romans existent, mais elles sont documentées et peuvent être intégrées à l'interprétation des résultats. Cette étape descriptive permet donc de mieux contrôler les effets liés à la longueur des œuvres avant l'application des méthodes de NMF et de K-means.

\subsubsection{Analyse des statistiques descriptives du nombre de mots par paquet}

\begin{table}[H]
\centering
\caption{Statistiques descriptives du nombre de mots par paquet}
\begin{tabular}{lr}
\toprule
\textbf{Indicateur} & \textbf{Valeur} \\
\midrule
Nombre de paquets & 5327 \\
Moyenne & 785.30 \\
Écart-type & 212.13 \\
Minimum & 42 \\
Premier quartile & 636 \\
Médiane & 753 \\
Troisième quartile & 902.5 \\
Maximum & 2944 \\
\bottomrule
\end{tabular}
\end{table}

Les statistiques descriptives du nombre de mots par paquet permettent d'évaluer la taille effective des unités textuelles utilisées pour l'analyse thématique. Le corpus segmenté comprend 5327 paquets, avec une moyenne d'environ 785 mots par paquet. La médiane, située à 753 mots, est relativement proche de cette moyenne, ce qui indique que la majorité des paquets se situe autour d'un volume textuel comparable.

L'écart-type est d'environ 212 mots, ce qui traduit une dispersion modérée autour de la moyenne. Le coefficient de variation est d'environ 27 \%, puisqu'il correspond au rapport entre l'écart-type et la moyenne. Cette valeur indique que les paquets ne sont pas parfaitement homogènes en longueur, mais que la variabilité reste acceptable pour une analyse thématique. Cette variation s'explique par le fait que les paquets sont construits à partir d'un nombre fixe de phrases, et non d'un nombre fixe de mots : selon le style, le rythme narratif ou la longueur syntaxique des phrases, deux paquets de 40 phrases peuvent contenir un nombre de mots assez différent.

Les quartiles permettent de préciser cette distribution. Un quart des paquets contient au maximum 636 mots, tandis que la moitié des paquets contient au maximum 753 mots. Le troisième quartile, situé à 902,5 mots, indique que 75 \% des paquets ne dépassent pas environ 903 mots. Ainsi, la majorité des unités d'analyse se concentre dans une zone relativement stable, comprise approximativement entre 636 et 903 mots pour la moitié centrale du corpus. Cette concentration est favorable au topic modeling, car elle limite les écarts de taille entre les unités comparées.

Les valeurs extrêmes doivent toutefois être prises en compte. Le paquet le plus court contient seulement 42 mots, tandis que le plus long atteint 2944 mots. Ces deux valeurs indiquent la présence d'unités atypiques. Le paquet de 42 mots correspond probablement à une fin de roman ou à un segment résiduel composé de moins de phrases que les autres. À l'inverse, le paquet de 2944 mots peut provenir d'un passage composé de phrases particulièrement longues ou d'une segmentation imparfaite du texte source. Ces cas extrêmes peuvent influencer l'analyse, car un paquet très long contient davantage d'information lexicale et peut donc peser plus fortement dans la représentation des thèmes.

Dans l'ensemble, ces statistiques montrent que la segmentation en paquets de 40 phrases produit des unités textuelles globalement comparables. Malgré une variabilité réelle, la distribution reste suffisamment concentrée autour de la médiane pour permettre une analyse thématique cohérente. Il serait toutefois pertinent d'examiner les paquets extrêmes afin de vérifier s'ils correspondent à des cas normaux du corpus ou à des problèmes de segmentation.


\subsection{Lemmatisation du corpus de Zola}

Après avoir obtenu un corpus segmenté en paquets de phrases et relativement homogène en termes de nombre de mots, l'étape suivante de préparation des données a consisté à effectuer une lemmatisation du corpus\footnote{Code disponible en annexe, dans la section « Lemmatisation et filtrage lexical ».}.

Pour réaliser cette lemmatisation, nous avons utilisé la bibliothèque \texttt{spaCy} en Python, qui offre des outils performants pour le traitement automatique du langage. La lemmatisation permet de regrouper les différentes formes d'un même mot sous une seule forme de base : par exemple, les mots « manger », « mange » et « manges » sont tous réduits au lemme « manger ». Cette étape facilite l'identification des thèmes et contribue à réduire la dimensionnalité du corpus.

Comme pour la segmentation en paquets de phrases, plusieurs essais ont été réalisés avec la bibliothèque de lemmatisation \texttt{spaCy}.

Nous avons choisi d'utiliser le modèle \texttt{fr\_core\_news\_lg} de \texttt{spaCy}, qui est un modèle de langue française de grande taille. Ce choix a été motivé par la richesse lexicale et la complexité syntaxique des romans de Zola, qui nécessitent un modèle capable de gérer une grande variété de formes linguistiques. Le modèle \texttt{fr\_core\_news\_lg} a été retenu, car il offre une lemmatisation plus précise que des modèles plus légers\footnote{Les modèles \texttt{fr\_core\_news\_sm} et \texttt{fr\_core\_news\_md} ont également été envisagés.}, ce qui est important pour l'identification des thèmes dans les textes littéraires.


Une liste de mots à exclure a été construite progressivement au cours de plusieurs itérations de prétraitement\footnote{Elle est présente en annexe du document.}. Elle résulte de l'observation conjointe des sorties de lemmatisation et de la matrice terme-document. L'objectif était d'identifier les formes lexicales très fréquentes dans le corpus, mais peu discriminantes pour l'analyse thématique.

Cette liste contient d'abord des mots-outils, tels que les pronoms personnels, les conjonctions, les déterminants et les prépositions. Ces termes structurent grammaticalement les phrases, mais ils ne permettent pas d'identifier les thèmes dominants des romans. Leur présence importante dans la matrice terme-document risquait donc de réduire la lisibilité des topics produits par la NMF.

La liste a également été enrichie par l'ajout de certains noms propres, notamment des noms de personnages récurrents dans les romans de Zola. Ces noms apparaissent fréquemment dans les textes, mais leur poids lexical tendait à orienter les topics vers des personnages plutôt que vers des thèmes plus généraux, comme le travail, la famille, la religion, l'argent ou les rapports sociaux. Leur exclusion permet donc de privilégier une interprétation thématique plutôt qu'une simple identification des protagonistes.

Ce filtrage relève d'un choix méthodologique empirique. Lors des premiers essais, plusieurs mots restaient fortement représentés dans les résultats malgré les paramètres appliqués au lemmatiseur et au vectoriseur. La liste d'exclusion a donc été ajustée à partir de l'examen des termes les plus fréquents et des premiers topics obtenus. Elle vise à améliorer la qualité interprétative de la matrice terme-document en conservant prioritairement les mots porteurs d'information thématique.

Une fonction applique le modèle \texttt{spaCy} à chaque segment du corpus afin d'en extraire une version lemmatisée et filtrée\footnote{Fonction en Python disponible en annexe du document.}. Pour chaque token, le lemme est d'abord converti en minuscules afin de normaliser les formes lexicales. La fonction exclut ensuite les mots vides, la ponctuation, les nombres, les espaces, ainsi que les lemmes trop courts ou présents dans la liste personnalisée de mots à retirer. Le filtrage conserve uniquement les noms et les adjectifs, car ces catégories grammaticales sont généralement les plus porteuses d'information thématique dans le cadre d'une analyse de topics. Le résultat final est une chaîne de caractères composée des lemmes retenus, qui peut ensuite être utilisée pour construire la matrice terme-document.

À l'issue de cette préparation, chaque roman est représenté par un ensemble de segments lemmatisés et filtrés. Ces segments constituent les unités d'analyse utilisées pour les deux méthodes de modélisation thématique présentées dans la section suivante.
